\begin{afterword}
\summaryhead

The post introduces ``thingspace'': a space with a dimension for every property, in which any object is a point. Colours become coordinates, and a robin's volume, mass, DNA, shape and colour all become parts of its position. Uncertainty about a robin becomes a cloud of probability.

The author then turns to radial categories. The central ``mother'' conceives a child, gives birth to it and raises it; genetic, surrogate and adoptive mothers are variants. The intension of such a word can be pictured as a glow in thingspace, brightest at the centre, and its extension as clusters of points: robins and sparrows at the centre of birdness, penguins and ostriches at the edge. The author holds that the glow was learned from the clusters. Real clusters need have no set of necessary and sufficient properties, so a simple definition will meet exceptions. That is acceptable. A definition such as ``feathered flying things'' only has to lead the listener to the cluster, and objecting to its exceptions is nitpicking.

\responsehead

The post is mainly a picture, and a useful one. It makes correlations between properties visible, and it gives later posts a vocabulary of clusters and glows. Its arithmetic and geometry are right as far as they go.

Its ideas about categories are older than the post, which does not claim otherwise. ``Radial categories'' and the ``mother'' example are George Lakoff's, from the book the previous post cited. That a natural group may share no single defining feature is Wittgenstein's family resemblance, which Rosch took up in the work on prototypes. A reader would not learn from this post where the ideas come from, and the notes give the sources.

The one claim the post makes in its own voice is that the category was learned from the clusters: ``First came the structure-in-the-world \ldots\ then, by observing it, we formed a category.'' The author offers it as a view (``as I see it''), and the picture cannot support it. The birds were placed with a distance ``that corresponds as well as possible to perceived similarity in humans,'' so the clusters are clusters in human judgment, and the picture cannot tell us whether the judgments followed the world or the categories shaped the judgments. Rosch, stating the same principle, said that it concerned ``the perceived world,'' and that the categories a culture already has may help decide which attributes people perceive. The post's order, world first and category after, may be largely right for birds; the post does not show it.

The practical conclusion is modest and reasonable. A definition that leads a listener to the right cluster has done its work. This restates, as a result of the new picture, what ``Similarity Clusters'' said two days earlier.

In short: a useful picture of categories as clusters, drawing on Lakoff, Wittgenstein and Rosch; its one claim of its own, that categories follow the clusters in the world, is offered as a view, and a picture whose distances are human judgments cannot show it.
\end{afterword}
